%==============================================================================
% Sjabloon poster bachproef
%==============================================================================
% Gebaseerd op document class `a0poster' door Gerlinde Kettl en Matthias Weiser
% Aangepast voor gebruik aan HOGENT door Jens Buysse en Bert Van Vreckem

\documentclass[a0,portrait]{hogent-poster}

% Info over de opleiding
\course{Bachelorproef}
\studyprogramme{toegepaste informatica}
\academicyear{2025-2026}
\institution{Hogeschool Gent, Valentin Vaerwyckweg 1, 9000 Gent}

% Info over de bachelorproef
\title{XR SkillSim: nieuwe medische XR-oefeningen bouwen in vier stappen}
\subtitle{Casus: BLS- en AED-training voor verpleegkunde op Meta Quest 3}
\author{Jesse Vaerendonck}
\email{jesse.vaerendonck@student.hogent.be}
\supervisor{Dhr. J. Van Herreweghe}
\cosupervisor{Dhr. T. Cuelenaere}

% Indien ingevuld, wordt deze informatie toegevoegd aan het einde van de
% abstract. Zet in commentaar als je dit niet wilt.
%\specialisation{Mobile \& Enterprise Development} % TODO: pas aan naar jouw afstudeerrichting
\keywords{Extended Reality, Unity, state machine, simulatietraining, Meta Quest 3}
%\projectrepo{https://github.com/user/repo} % TODO: vul je repo-URL in of laat in commentaar

\begin{document}

\maketitle

\begin{abstract}
Zorgopleidingen willen vaardigheden schaalbaar en herhaalbaar laten oefenen, maar XR-oefeningen worden vaak ad hoc gebouwd en zijn daardoor moeilijk uit te breiden. Deze bachelorproef toont hoe een modulaire architectuur het mogelijk maakt om medische XR-trainingsoefeningen systematisch op te bouwen, met objectieve logdata als basis voor feedback. Als bewijs werd een volledige reanimatieoefening (BLS en AED) gerealiseerd op Meta Quest 3.
\end{abstract}

\begin{multicols}{2} % This is how many columns your poster will be broken into, a portrait poster is generally split into 2 columns

\section{Inleiding}

Verpleegkundestudenten moeten procedures zoals reanimatie veilig, herhaaldelijk en scenario-gericht kunnen inoefenen. Simulatielaboratoria zijn echter logistiek beperkt: elk oefenmoment vraagt planning, infrastructuur en docentcapaciteit, waardoor studenten minder vaak kunnen oefenen op besluitvorming onder tijdsdruk. Extended Reality maakt herhaalbaar en schaalbaar oefenen mogelijk, maar in de praktijk wordt elke XR-oefening vaak ad hoc en vanaf nul gebouwd. Dat is duur, foutgevoelig en moeilijk uit te breiden naar nieuwe scenario's.

Deze bachelorproef verlegt daarom de focus van één oefening naar de onderliggende bouwstenen: hoe ontwerp je een XR-omgeving waarin nieuwe medische trainingsscenario's systematisch aangemaakt, gevalideerd en gelogd kunnen worden? De centrale onderzoeksvraag luidt: in hoeverre biedt zo'n XR-applicatie een veilige en bruikbare leeromgeving, waarbij objectieve data over de uitvoering inzetbaar zijn voor feedback en evaluatie? Een richtlijnconforme reanimatieoefening (BLS en AED) dient als casus om de aanpak aan te tonen.

\section{Methode}

Het onderzoek volgt Design Science Research: een artefact wordt iteratief ontworpen, ontwikkeld en geëvalueerd in functie van een concreet praktijkprobleem. Uit literatuur en de ERC-richtlijnen werden toetsbare requirements afgeleid, vertaald naar een ontwerp en gerealiseerd als proof of concept.

De technische realisatie gebeurde in Unity 6 met OpenXR, de Meta XR Core SDK, de Meta Interaction SDK en de XR Interaction Toolkit, standalone op de Meta Quest 3. De evaluatie bestond uit interne technische verificatie en scenario-walkthroughs, waarbij de uitgevoerde handelingen werden vergeleken met de geregistreerde logdata.

\section{Resultaten}

De architectuur steunt op drie herbruikbare lagen: een state machine die elk scenario opdeelt in procedurele stappen, event-driven communicatie die componenten ontkoppelt, en een loggingsysteem (xAPI-geïnspireerd, JSON met tijdstempels) dat elke interactie automatisch registreert.

\begin{center}
  \captionsetup{type=figure}
  \includegraphics[width=1.0\linewidth]{architecture}
  \captionof{figure}{Modulaire systeemarchitectuur met event-driven communicatie.}
\end{center}

Daarbovenop volgt het aanmaken van een nieuwe oefening steeds hetzelfde stramien van vier stappen:

\begin{enumerate}
  \item de procedure opdelen in scenario-states met overgangsvoorwaarden;
  \item interactiezones (colliders) en interactors koppelen aan de relevante objecten;
  \item validatieregels definiëren per stap (volgorde, timing, plaatsing);
  \item logging gebeurt automatisch via de event-laag.
\end{enumerate}

\begin{center}
  \captionsetup{type=figure}
  \includegraphics[width=1.0\linewidth]{ScenarioFlowState}
  \captionof{figure}{Scenarioflow: elke stap in de procedure is een afzonderlijke state.}
\end{center}

Met dit stramien werd de casus gebouwd: een richtlijnconforme reanimatieoefening met alarmeren, compressies, AED aanbrengen en veilig schokken. Herbruikbare componenten zoals een velocity-based CPRInteractor (compressiedetectie) en een ruimtelijk validatiesysteem voor AED-elektroden werden toegevoegd. In de walkthroughs kwam de gelogde eventvolgorde telkens overeen met de uitgevoerde handelingen, zodat volgorde, timing en fouten objectief analyseerbaar zijn.

\section{Discussie en besluit}

De meerwaarde van dit onderzoek is dat XR-trainingsoefeningen niet langer ad hoc gebouwd hoeven te worden: de modulaire architectuur maakt nieuwe scenario's systematisch aanmaakbaar, met objectieve procesdata als automatisch bijproduct voor feedback en evaluatie. De BLS-casus bewijst de technische en procedurele haalbaarheid op betaalbare standalone hardware.

Zonder fysieke tegendruk blijft de applicatie wel een procedureel en cognitief oefeninstrument, geen vervanger van fysieke vaardigheidstraining. Brede gebruikerstests met studenten en externe experten vormen de volgende stap, samen met multi-user oefeningen, hybride simulatie met fysieke manikins en dashboarding voor docenten. Zo groeit XR SkillSim uit tot fundament voor de HOGENT XR Academy.

\section{Referenties}

{\small
Peffers, K. e.a. (2007). A Design Science Research Methodology for Information Systems Research. \textit{Journal of MIS}, 24(3).

European Resuscitation Council (2025). \textit{ERC Guidelines: Basic Life Support}.

Greif, R. e.a. (2023). ILCOR: Immersive technologies in resuscitation education.
}

\end{multicols}
\end{document}
